\chapter{Sprint 6 : Collaboration, restitution et industrialisation}
\label{chap:sprint6}

\sectionsansnum{Introduction}

Ce dernier sprint couvre trois volets complémentaires : le travail à plusieurs, la production
du document qui sort du système, et l'industrialisation du déploiement.

\section{Objectif du sprint}

L'interprétation d'un enregistrement difficile se fait rarement seul. Le premier objectif
était donc de permettre à un praticien de \textbf{solliciter un confrère}, en gardant la
discussion attachée à l'examen concerné.

Le deuxième objectif était la \textbf{restitution}. Un système d'aide au diagnostic dont on ne
peut rien extraire n'a pas d'usage : il fallait produire un compte rendu composable, adapté au
destinataire.

Le troisième était l'\textbf{industrialisation} : rendre le déploiement reproductible et
vérifier automatiquement qu'une modification ne casse rien.

\section{Analyse du sprint 6}

\subsection{Sprint Backlog}

\begin{table}[H]
\centering\small
\renewcommand{\arraystretch}{1.15}
\begin{tabular}{@{}c|L{2.5cm}|L{4.6cm}|L{4.6cm}|c@{}}
\hline
\ligneentete \entete{ID} & \entete{Module} & \entete{R\'ecit utilisateur} & \entete{T\^aches} & \entete{Estim.} \\
\hline
\multirow{3}{*}{US-47\,\`a\,49} & \multirow{3}{=}{Collaboration} & \multirow{3}{=}{En tant que m\'edecin, je veux solliciter un confr\`ere sur un examen pr\'ecis} & Mod\'eliser fils et messages rattach\'es \`a un examen & \multirow{3}{*}{18 pts} \\
 &  &  & Cr\'eer l'interface de discussion &  \\
 &  &  & Centraliser les sollicitations re\c{c}ues &  \\
\hline
\multirow{3}{*}{US-45, 46, 55} & \multirow{3}{=}{Annotation et normes} & \multirow{3}{=}{En tant que m\'edecin, je veux annoter une \'epoque et situer chaque m\'etrique par rapport aux normes} & Enregistrer les notes horodat\'ees \`a l'\'epoque & \multirow{3}{*}{15 pts} \\
 &  &  & Synchroniser l'hypnogramme annot\'e &  \\
 &  &  & Comparer les m\'etriques aux valeurs de r\'ef\'erence &  \\
\hline
\multirow{3}{*}{US-21, 54} & \multirow{3}{=}{Restitution} & \multirow{3}{=}{En tant que m\'edecin, je veux composer et exporter mon compte rendu} & S\'electionner les sections \`a inclure & \multirow{3}{*}{13 pts} \\
 &  &  & Produire le document pagin\'e &  \\
 &  &  & Afficher le tableau de bord de la patient\`ele &  \\
\hline
\multirow{4}{*}{US-56\,\`a\,60} & \multirow{4}{=}{Industrialisation} & \multirow{4}{=}{En tant que d\'eveloppeur, je veux un d\'eploiement reproductible et v\'erifi\'e automatiquement} & Conteneuriser les trois services & \multirow{4}{*}{37 pts} \\
 &  &  & Orchestrer la pile en une commande &  \\
 &  &  & Automatiser contr\^oles et construction d'images &  \\
 &  &  & Traduire l'interface et ajouter le th\`eme sombre &  \\
\hline
\end{tabular}
\caption{Sprint Backlog --- Sprint 6}
\label{tab:bl-s6}
\end{table}

\subsection{Diagramme des cas d'utilisation}

\begin{figure}[H]
\centering
\begin{tikzpicture}[
    x=1cm,y=1cm,
    uc/.style     = {ellipse, draw=hypnoblue!60, fill=hypnolight,
                     align=center, font=\scriptsize, text width=2.7cm,
                     inner sep=1.5pt, minimum height=1.0cm},
    ucp/.style    = {ellipse, draw=hypnored!70, fill=hypnored!7,
                     align=center, font=\scriptsize\bfseries, text width=2.7cm,
                     inner sep=1.5pt, minimum height=1.0cm},
    lien/.style   = {hypnogray!75, line width=0.55pt},
    rel/.style    = {-{Latex[length=2mm]}, hypnoblue!70, line width=0.5pt, dashed},
    et/.style     = {font=\tiny\itshape, text=hypnoblue!85, fill=white, inner sep=1pt},
    nomact/.style = {font=\scriptsize\bfseries, align=center, text width=2.2cm},
  ]

  \tikzset{acteur/.pic={
      \draw[hypnoblue,line width=0.8pt] (0,0) circle (0.17);
      \draw[hypnoblue,line width=0.8pt] (0,-0.17) -- (0,-0.66);
      \draw[hypnoblue,line width=0.8pt] (-0.28,-0.34) -- (0.28,-0.34);
      \draw[hypnoblue,line width=0.8pt] (0,-0.66) -- (-0.24,-1.04);
      \draw[hypnoblue,line width=0.8pt] (0,-0.66) -- (0.24,-1.04);}}

  \draw[hypnoblue!55,line width=0.9pt,rounded corners=3pt]
       (3.9,-4.6) rectangle (12.0,2.4);
  \node[font=\small\bfseries,text=hypnoblue,anchor=north west] at (4.15,2.25)
       {Collaboration et restitution};

  \node[ucp] (col) at (6.2, 1.2)  {Solliciter un\\confrère};
  \node[uc]  (msg) at (10.0, 1.2) {Échanger des\\messages};
  \node[uc]  (rec) at (10.0,-0.4) {Consulter la boîte\\de réception};
  \node[uc]  (ann) at (6.2,-1.1)  {Annoter une\\époque};
  \node[uc]  (nor) at (10.0,-2.0) {Comparer aux\\normes};
  \node[ucp] (pdf) at (6.2,-3.3)  {Composer le\\compte rendu};
  \node[uc]  (sec) at (10.0,-3.6) {Choisir les\\sections};

  \pic at (1.6,0.7) {acteur};
  \node[nomact,anchor=north] at (1.6,0.05) {Médecin};

  \pic at (14.1,0.7) {acteur};
  \node[nomact,anchor=north] at (14.1,0.05) {Confrère};

  \foreach \c in {col,ann,pdf} \draw[lien] (2.0,0.35) -- (\c.west);
  \draw[lien] (13.7,0.35) -- (msg.east);
  \draw[lien] (13.7,0.35) -- (rec.east);

  \draw[rel] (col) -- node[et,pos=0.5] {$\ll$include$\gg$} (msg);
  \draw[rel] (col) -- node[et,pos=0.6] {$\ll$extend$\gg$} (rec);
  \draw[rel] (ann) -- node[et,pos=0.6] {$\ll$extend$\gg$} (nor);
  \draw[rel] (pdf) -- node[et,pos=0.55] {$\ll$include$\gg$} (sec);

\end{tikzpicture}
\caption{Diagramme des cas d'utilisation du sprint 6 --- collaboration et restitution}
\label{fig:uc-s6}
\end{figure}

\subsubsection{Description textuelle des cas d'utilisation}

\begin{casutilisation}{Description textuelle du cas « Solliciter un confrère »}
Titre & Solliciter un confrère sur un examen \\ \hline
Acteurs & Médecin demandeur, confrère sollicité \\ \hline
Description & Ouvrir une discussion rattachée à un examen afin de recueillir un avis, en
conservant le contexte clinique \\ \hline
Préconditions & Les deux praticiens sont authentifiés ; l'examen existe \\ \hline
Scénario de base &
\begin{minipage}[t]{\linewidth}
\begin{enumerate}[leftmargin=1.1em,itemsep=1pt,topsep=3pt]
  \item Le médecin ouvre la fiche du patient et sélectionne l'examen.
  \item Il demande la liste des confrères disponibles.
  \item Il sélectionne un confrère et rédige sa demande.
  \item Le système crée le fil, ou réutilise celui qui existe déjà, et enregistre le message.
  \item Le confrère consulte sa boîte de réception et répond dans le même fil.
\end{enumerate}
\end{minipage} \\ \hline
Scénario alternatif &
\begin{minipage}[t]{\linewidth}
\begin{itemize}[leftmargin=1.1em,itemsep=1pt,topsep=3pt]
  \item \textbf{4a} --- Un fil existe déjà entre ces deux praticiens sur cet examen : le
        message y est ajouté plutôt que d'ouvrir un doublon.
\end{itemize}
\end{minipage}
\end{casutilisation}

\section{Réalisation}

\subsection{Consultations confraternelles}

Un fil de discussion est rattaché à un examen et à exactement deux praticiens. L'intérêt de ce
dispositif est que la discussion reste attachée à l'examen : le correspondant dispose du
contexte sans qu'il faille le lui transmettre. Une vue centralisée rassemble toutes les
sollicitations reçues.

\capturefig{consultation}{0.86}{Fil de discussion rattaché à un examen}{fig:consultation}

\subsection{Annotation et comparaison aux normes}

Le praticien peut annoter une époque précise de l'hypnogramme ; ses observations sont
conservées et peuvent être synchronisées vers le dossier, ce qui archive l'hypnogramme annoté
aux côtés des autres pièces.

Chaque métrique AASM est par ailleurs accompagnée de sa plage de référence et d'une comparaison
aux valeurs de la cohorte. Une valeur hors norme est signalée visuellement, et une liste
d'alertes cliniques reprend en clair les anomalies détectées.

\subsection{Production du compte rendu}

Le compte rendu est produit côté client : les sections retenues sont converties en images par
capture du rendu, puis assemblées dans un document paginé. Cette approche garantit que le
document reproduit exactement ce que le praticien a validé à l'écran, sans qu'un moteur de
rendu distinct puisse introduire un écart.

Le praticien choisit les sections à inclure : hypnogramme, métriques AASM, classe de sévérité,
facteurs explicatifs, tracés bruts et comparaison aux normes. Cette composition à la carte
répond à des destinataires différents --- un confrère n'attend pas le même document qu'un
dossier d'assurance. Le document est ensuite archivé au même titre que les autres pièces.

\capturefig{rapport}{0.78}{Compte rendu clinique généré}{fig:rapport}

\subsection{Multilinguisme et thème sombre}

L'interface est intégralement traduite en français et en anglais, les chaînes étant centralisées
dans un catalogue unique et la bascule immédiate. Un thème sombre est également disponible,
adapté à la lecture nocturne --- un praticien qui interprète un enregistrement le fait rarement
en plein jour.

\capturefig{dark}{0.86}{Tableau de bord du médecin en thème sombre}{fig:dark}

\subsection{Conteneurisation}

La pile est décrite par un fichier d'orchestration unique qui déclare trois services.

\begin{table}[H]
\centering

\small
\begin{tabular}{@{}lllL{5.4cm}@{}}
\toprule
\textbf{Service} & \textbf{Image de base} & \textbf{Port} & \textbf{Particularité} \\
\midrule
Base de données & \texttt{postgres:15-alpine} & 5432 &
Volume persistant, sonde de disponibilité \\
\addlinespace[2pt]
Interface applicative & \texttt{python:3.10-slim} & 8000 &
PyTorch en version processeur installé en premier \\
\addlinespace[2pt]
Serveur web & \texttt{node:20-alpine} puis \texttt{nginx:alpine} & 80 &
Construction en deux étapes, redirection du routage client \\
\bottomrule
\end{tabular}
\caption{Services de la pile conteneurisée}
\label{tab:services}
\end{table}

Le fichier de construction de l'interface applicative comporte une particularité qui mérite
explication : PyTorch y est installé depuis l'index dédié aux processeurs, \emph{avant} les
autres dépendances. Sans cette précaution, l'installation télécharge plusieurs gigaoctets de
bibliothèques inutiles en déploiement --- l'inférence en production s'exécutant sur processeur
--- et fait échouer la construction sur une machine de développement.

La sonde de disponibilité de la base est également indispensable : sans elle, l'interface
applicative démarrerait avant que la base n'accepte les connexions et échouerait à créer son
schéma.

\subsection{Intégration continue}

Une chaîne d'intégration continue s'exécute à chaque livraison sur la branche principale et à
chaque proposition de fusion.

\begin{table}[H]
\centering

\small
\begin{tabular}{@{}lL{10.4cm}@{}}
\toprule
\textbf{Travail} & \textbf{Étapes} \\
\midrule
Interface applicative &
Installation des dépendances avec PyTorch en version processeur ; contrôle syntaxique
de tous les modules ; démarrage d'une instance PostgreSQL de service ; démarrage réel
de l'application ; vérification du schéma OpenAPI et de la présence des points
d'entrée attendus ; contrôle du cloisonnement des rôles \\
\addlinespace[3pt]
Interface utilisateur &
Installation reproductible des dépendances ; analyse statique ; construction de
production ; publication du bundle en tant qu'artefact \\
\addlinespace[3pt]
Images de conteneur &
Construction des deux images avec cache de couches, conditionnée à la réussite des
deux travaux précédents \\
\bottomrule
\end{tabular}
\caption{Travaux de la chaîne d'intégration continue}
\label{tab:ci}
\end{table}

Le travail portant sur l'interface applicative ne se contente pas d'un contrôle syntaxique : il
\textbf{démarre réellement l'application} contre une base éphémère, puis exerce trois
vérifications de comportement. Un appel sans jeton doit recevoir un refus d'authentification ;
une connexion administrateur doit délivrer un jeton ; et ce jeton d'administrateur, présenté au
registre des patients, doit recevoir un refus d'autorisation. Ces trois assertions couvrent
l'essentiel du cloisonnement des rôles et échoueraient immédiatement si une dépendance de
sécurité venait à disparaître au fil d'une modification.

\capturefig{ci}{0.9}{Exécution de la chaîne d'intégration continue}{fig:ci}

\subsection{Validation fonctionnelle}

En complément, une campagne de scénarios fonctionnels a été exécutée manuellement à chaque fin
de sprint.

\begin{table}[H]
\centering

\small
\begin{tabular}{@{}clL{7.2cm}c@{}}
\toprule
\textbf{N\textsuperscript{o}} & \textbf{Scénario} & \textbf{Critère de réussite} &
\textbf{État} \\
\midrule
1 & Connexion valide & Jeton délivré, espace de travail ouvert & \checkmark \\
2 & Connexion refusée & Compte suspendu ou licence expirée rejeté avec message
    distinct & \checkmark \\
3 & Cloisonnement des rôles & Un médecin n'atteint aucune route d'administration &
    \checkmark \\
4 & Cloisonnement des patients & Un médecin n'accède pas au dossier d'un confrère &
    \checkmark \\
5 & Analyse complète & Enregistrement déposé, hypnogramme et métriques affichés &
    \checkmark \\
6 & Montage incomplet & Dérivation absente : repli appliqué, analyse aboutie &
    \checkmark \\
7 & Fichier invalide & Format non reconnu refusé avec message explicite & \checkmark \\
8 & Prédiction de sévérité & Classe, probabilités et facteurs restitués & \checkmark \\
9 & Import de rapport externe & Formulaire prérempli depuis un fichier tabulaire ou
    balisé & \checkmark \\
10 & Archivage en arrière-plan & Résultats affichés avant la fin du transfert &
    \checkmark \\
11 & Échec d'archivage & Transfert interrompu : examen consultable, erreur signalée &
    \checkmark \\
12 & Lien expiré & Lien réutilisé après échéance : accès refusé & \checkmark \\
13 & Annotation d'époque & Note enregistrée et retrouvée au rechargement & \checkmark \\
14 & Consultation confraternelle & Message émis par un praticien, reçu par l'autre &
    \checkmark \\
15 & Compte rendu & Document produit avec les seules sections sélectionnées &
    \checkmark \\
16 & Journal d'audit & Actions sensibles tracées et exportées & \checkmark \\
17 & Bascule linguistique & Aucune chaîne non traduite sur le parcours principal &
    \checkmark \\
18 & Démarrage de la pile & Trois services opérationnels en une commande & \checkmark \\
\bottomrule
\end{tabular}
\caption{Scénarios de validation fonctionnelle}
\label{tab:tests}
\end{table}

Il faut reconnaître la limite de cette approche : une campagne manuelle n'est ni reproductible
ni exhaustive, et elle ne protège pas contre les régressions entre deux exécutions. Les trois
assertions automatisées constituent un premier filet de sécurité, mais une véritable couverture
de tests reste à construire.

\sectionsansnum{Conclusion}

Ce dernier sprint a complété la plateforme sur ses trois derniers volets. Les praticiens peuvent
solliciter un confrère sans quitter le contexte de l'examen, composer un compte rendu adapté à
son destinataire, et travailler dans la langue et le thème de leur choix.

La pile est entièrement conteneurisée et vérifiée par une chaîne d'intégration continue qui
démarre réellement l'application et contrôle le cloisonnement des rôles.

Le projet est ainsi complet sur son périmètre fonctionnel. La conclusion générale en dresse le
bilan et ouvre sur les perspectives.
